% !TEX program = lualatex
\documentclass[a4paper,11pt]{scrartcl}

\usepackage[
  lang=zh,
  mode=journal,
  theme=graphite,
  fontset=lm,
  toc=true,
  toc-layout=page
]{synaptic}
\usepackage{booktabs,tabularx}
\newcolumntype{Y}{>{\raggedright\arraybackslash}X}

\SynapticHeader{synaptic 技术参考}
\SynapticSubHeader{架构、不变量与发布流程}
\SynapticPreprint{v2.2.0}
\SynapticTitle{synaptic 技术参考}
\SynapticAuthor{synaptic 项目}
\SynapticAffiliation{开源 LaTeX 社区}
\SynapticSubject{synaptic 2.2.0 技术参考}
\SynapticKeywords{LaTeX3, 架构, l3build, 回归测试}

\begin{abstract}
  本文档说明 synaptic 的模块图、加载期不变量、命名空间规则、字体与颜色子系统、
  扩展点、生成文件策略以及自动化验证流程。
\end{abstract}

\begin{document}
\SynapticMakeTitle

\section{架构}

\begin{verbatim}
synaptic.sty
  `-- synaptic-base.sty
      |-- synaptic-lang-{en,zh}.def
      |-- synaptic-color.sty
      |-- synaptic-fonts.sty
      |-- synaptic-layout.sty
      |-- synaptic-theorem.sty
      |-- synaptic-title.sty       公共元数据与渲染器
      `-- 模式模块
          |-- synaptic-journal.sty     journal
          |-- synaptic-book.sty        book
          |-- synaptic-lecture.sty     lecture
          |-- synaptic-boxes.sty       lecture
          `-- synaptic-notes.sty       notes
\end{verbatim}

薄封装层把用户选项转发给 \texttt{synaptic-base}，并定义少量不易冲突的便捷命令。
基础模块负责检查引擎与文档类、解析选项、加载标签、验证书籍模式，然后分发核心
模块和模式模块。

\noindent\begin{tabularx}{\textwidth}{@{}l l Y@{}}
  \toprule
  \textbf{模式} & \textbf{文档类不变量} & \textbf{附加模块} \\
  \midrule
  journal & KOMA 文章/报告类 & journal \\
  book & 含章和前文接口的 KOMA 类 & book \\
  lecture & KOMA 文章/报告类 & lecture, boxes \\
  notes & KOMA 文章/报告类 & notes \\
  \bottomrule
\end{tabularx}

\section{配置生命周期}

宏包级 \texttt{synaptic} 键系统在模块加载前执行。模式、语言、字体组、标题布局、
版心尺度、密度、编号、输出颜色模式、数学扩展和断行策略因此成为结构不变量。
运行时键系统只暴露无需重载宏包就能应用的
\texttt{theme}、\texttt{toc} 和 \texttt{toc-layout}。冻结键与未知键会产生明确警告。

书籍模式要求章和前文接口同时存在。LuaLaTeX 与 KOMA-Script 要求在不满足时立即给出
致命诊断，而不是延迟到后续的未定义命令。

\section{子系统}

\subsection{颜色令牌}

内置主题登记在全局属性表中。启用主题时，主色、辅色与强调色会映射到公开语义名，
并派生表面、边框、分隔线、链接、可读弱化文字、装饰浅色与正文颜色。print 与 mono
输出模式只改变派生角色，不改变组件代码。危险、警告和成功色是稳定的语义常量，不随
品牌主色改变含义。

\begin{verbatim}
\SynapticDefineTheme{名称}{主色}{辅色}[强调色]
\SynapticUseTheme{名称}
\end{verbatim}

参数是 xcolor 表达式。定义自定义主题时会同时登记并启用它。

\subsection{字体选择}

字体发现完全由 \verb|\IfFontExistsTF| 实现。每个候选项都是完整的衬线/无衬线/等宽/数学
组合；自动选择不会启用缺少配套字体的家族。固定优先级为 XCharter、Libertinus、
STIX Two、Latin Modern。显式选择不完整字体组时，报错会列出必需字体。

\texttt{lang=zh} 会以无预设字体方案加载 ctex。Noto CJK SC 优先，Fandol 作为 TeX Live
后备。仅在优先字体没有原生字形时，才使用受控的伪粗体或伪斜体。

\subsection{版式与文档结构}

页面布局结合模式默认值、版心尺度与密度档位。journal 与 lecture 优先适中行长，notes
支持紧凑单栏/双栏，book 使用镜像内外边、装订补偿、外侧页眉和空白偶数页。孤行惩罚、
垂直节奏、列表间距、行距与段首缩进都是明确配置，
不再用全局方式隐藏坏盒子。

KOMA 标题字体与间距在 \texttt{begindocument/before} 钩子中安装。hyperref、bookmark、
microtype、图表标题标点、列表与模式化公式编号由布局模块统一配置。

\subsection{标题与元数据}

标题值保存在 expl3 词元表和序列中。作者与机构标记使用平行序列，机构与邮箱可重复。
公共元数据分别驱动期刊内联标题、书籍扉页序列、课程信息条和紧凑笔记信息条。标题命令
验证非空标题并控制可选目录位置，不再隐式重置用户页码。标准 LaTeX 元数据只是输入适配层，
不会重定义 \verb|\maketitle|。

PDF 信息以源词元传给 hyperref，由 hyperref 只执行一次 Unicode PDF 字符串转换。预先
\verb|\pdfstringdef| 再传给 \verb|\hypersetup| 会导致双重编码，因此已禁止这种路径。

\subsection{定理与盒子}

thmtools 在 amsthm 之上定义 plain、definition 和 remark 三种样式。内置定理家族
使用模式化共享计数器：书籍默认按章，其他模式默认按节。讲义模式用 tcolorbox 包装核心
定理环境；例题与练习可共享定理计数器，警告与注记默认不编号。新公开短名使用
\texttt{syn} 前缀；v2.0 通用名是有条件的兼容层。

\section{命名空间规则}

\noindent\begin{tabularx}{\textwidth}{@{}l l Y@{}}
  \toprule
  \textbf{范围} & \textbf{模式} & \textbf{示例} \\
  \midrule
  公开命令 & \texttt{\textbackslash Synaptic\ldots} & \SynapticCmd{SynapticDefineTheme} \\
  公开环境 & \texttt{syn...} & \texttt{synwarning} \\
  内部函数 & \texttt{\textbackslash\_\_synaptic\_\ldots:} & 标题渲染辅助函数 \\
  模块接口 & \texttt{\textbackslash synaptic\_\ldots:} & 模式访问器 \\
  局部状态 & \texttt{\textbackslash l\_synaptic\_\ldots} & 当前主题词元表 \\
  全局登记表 & \texttt{\textbackslash g\_synaptic\_\ldots} & 主题属性表 \\
  标签常量 & \texttt{\textbackslash c\_synaptic\_label\_\ldots} & 定理标签 \\
  \bottomrule
\end{tabularx}

新公开名应避免常见名词。兼容别名必须使用 \verb|\ProvideDocumentCommand| 或显式存在性检查。

\section{扩展点}

\begin{itemize}
  \item 在新的 \texttt{synaptic-lang-*.def} 中增加标签，再扩展语言键与定义文件分发。
  \item 新字体组需同时增加完整性谓词、设置函数、显式分支，并明确其在自动链中的位置。
  \item 新模式需扩展基础键系统、布局分发和模块加载器，并声明、测试必需的 KOMA 能力。
  \item 视觉扩展应基于语义颜色令牌，而不是固定 RGB 值。
\end{itemize}

\section{源码与发布流程}

\texttt{synaptic.dtx} 是唯一源头。DocStrip 通过 \texttt{synaptic.ins} 生成可安装的
\texttt{.sty} 和 \texttt{.def} 文件。\texttt{tex/latex/synaptic/} 下的生成副本为了零安装使用而
提交，必须与新解包结果字节级一致。

\begin{verbatim}
l3build unpack
l3build check
l3build doc
l3build ctan
\end{verbatim}

九组回归测试覆盖论文标题与声明、书籍前后文和章级编号、讲义导航与盒子、笔记卡和
标准元数据、结构选项、自定义主题与中文排版。CI 同时构建源码文档、四份手册和所有示例。

\section{兼容性说明}

2.2 保留小写书籍前后文命令和讲义盒子通用名，但仅作为有条件的 v2.0 兼容别名。
新文档应使用带命名空间的公开形式。已移除的 \texttt{style=} 和 \texttt{color=} 选项仍会报错，
并提供迁移提示。

\end{document}
